% Primera lectura: coordenadas homogéneas, la misma idea en cada dimensión, del álgebra a la geometría, transformaciones proyectivas
\section{Coordenadas homogéneas: puntos como rayos}\label{cap:05}\label{sec:homogeneas-primera}

Las fotografías de clase muestran el plano afín como el plano \(w=1\) dentro de un espacio
vectorial tridimensional. El punto afín \((x,y)\) se representa mediante
\[
  (x,y)\longleftrightarrow
  \begin{pmatrix}x\\y\\1\end{pmatrix}.
\]
Sin embargo, en geometría proyectiva no importa el representante concreto, sino el rayo que
parte del origen:
\[
\begin{pmatrix}x\\y\\1\end{pmatrix}
\sim
\lambda
\begin{pmatrix}x\\y\\1\end{pmatrix},
\qquad \lambda\neq0.
\]
Todos esos vectores pertenecen al mismo subespacio vectorial de dimensión uno y cortan al
plano \(w=1\) en el mismo punto afín.

\begin{center}
\begin{tikzpicture}[scale=.92]
  \coordinate (o) at (1,1);
  \draw[axisline,-{Latex[length=2.4mm]}] (o)--(8.2,1) node[right] {$y$};
  \draw[axisline,-{Latex[length=2.4mm]}] (o)--(-.2,-.1) node[below left] {$x$};
  \draw[axisline,-{Latex[length=2.4mm]}] (o)--(1,6) node[above] {$w$};
  \draw[teallight,fill=teallight!55,line width=.8pt]
    (-.2,3.2)--(6.8,3.2)--(8.2,4.25)--(1.2,4.25)--cycle;
  \node[labelback,text=tealdark] at (7.1,4.55) {plano \(w=1\)};
  \coordinate (p) at (4.8,3.74);
  \fill[rust] (p) circle (2.7pt) node[labelback,above] {\((x,y,1)\)};
  \draw[rust,line width=1.05pt,-{Latex[length=2.5mm]}] (o)--($(o)!1.38!(p)$);
  \fill[ink] (o) circle (2.4pt) node[labelback,below] {origen};
  \node[labelback,font=\small,text=rust] at (3.4,2.1)
    {un rayo representa un punto};
\end{tikzpicture}
\end{center}

\paragraph{Puntos en el infinito.}
Una recta afín con dirección \(v\) puede escribirse
\[
  p(\mu)=p+\mu v.
\]
Homogéneamente:
\[
\begin{pmatrix}p+\mu v\\1\end{pmatrix}
=
\begin{pmatrix}p\\1\end{pmatrix}
+\mu
\begin{pmatrix}v\\0\end{pmatrix}.
\]
Cuando \(|\mu|\) crece, la dirección dominante es \((v,0)\). La última coordenada nula
identifica un punto ideal o punto en el infinito. Todas las rectas paralelas con dirección
\(v\) comparten ese mismo punto ideal.

\paragraph{Primera traducción algebraica.}
Una matriz proyectiva general produce
\[
\begin{pmatrix}X\\Y\\W\end{pmatrix}
=H
\begin{pmatrix}x\\y\\1\end{pmatrix}.
\]
Si \(W\neq0\), se vuelve al plano afín normalizando:
\[
  (x',y')=\left(\frac{X}{W},\frac{Y}{W}\right).
\]
Las afinidades son el caso especial en que la última fila de \(H\) es \((0,0,1)\), por lo
que \(W\) permanece igual a \(1\).

\begin{idea}[En palabras del profesor]
``Terminamos con dos formas de expresión algebraica, una en \(\R^2\) y otra en el plano
proyectivo. [\ldots] Sobre el espacio proyectivo nos preguntamos especialmente qué pasaba con
los puntos impropios y cómo se expresaría una transformación con puntos fijos'' (resumen de clase 3, cursada 2025).
\end{idea}

\section{La misma idea en cada dimensión}\label{sec:misma-idea}

\textbf{Desarrollo.} Casi todo lo que la materia hace en \(\R^2\), \(\R^3\) y \(\R^n\) ya
existe en \(\R\), donde es tan simple que ni se nota. Leer cada concepto primero en la
recta, después en el plano y después en el espacio es la versión más chica de ``ir a lo
chico para volver a lo grande'': si la idea se entiende con un número, la matriz sólo
agrega direcciones. La clase 10 del profesor hace exactamente esto con la inversa: ``la
inversa es un elemento de un conjunto que, al operar con otro, resulta en el elemento
neutro respecto de esa operación'', sea la suma, el producto, la composición o el producto
de matrices.

{\small
\begin{tabularx}{\textwidth}{>{\bfseries\raggedright\arraybackslash}p{17mm}>{\raggedright\arraybackslash}p{26mm}>{\raggedright\arraybackslash}p{32mm}>{\raggedright\arraybackslash}p{32mm}Y}
\toprule
& \(\R\) & \(\R^2\) & \(\R^3\) y \(\R^n\) & Hecho geométrico\\
\midrule
Identidad & el número \(1\): \(1\cdot a=a\) & la matriz \(I\): \(Iv=v\) & \(I\) de
\(n\times n\) & la transformación que no mueve nada; es el neutro de la composición\\
\addlinespace
Inversa & \(1/a\); no existe si \(a=0\) & \(A^{-1}\); no existe si \(\det A=0\) & igual;
si \(\det A=0\) hay núcleo & deshacer el movimiento; imposible si alguna dirección fue
aplastada a un punto (clase 10)\\
\addlinespace
Tamaño con signo & \(|a|\) estira; el signo dice si da vuelta & \(\det A\): área orientada
del paralelogramo imagen del cuadrado unidad & volumen orientado; en \(\R^n\),
hipervolumen & cuánto se agranda una figura y si cambia de orientación\\
\addlinespace
Producto & \(a\cdot b\) & producto interno \(u\cdot v=|u||v|\cos\theta\) & igual; en
homogéneas, \(\ell\cdot p=0\) es incidencia & proyección de un vector sobre otro; cero si
son perpendiculares\\
\addlinespace
Reflexión y rotación & \(x\mapsto -x\): la única ``rotación'' de la recta es media vuelta
& multiplicar por \(e^{i\theta}\) en \(\mathbb{C}=\R^2\); reflejar es conjugar &
matriz de rotación \(3\times3\) con un eje; reflexión en un plano & conservar distancias;
la rotación conserva la orientación, la reflexión la invierte\\
\addlinespace
Autovector & \(a\,x=\lambda x\): trivial, todo \(x\) sirve con \(\lambda=a\) &
\(Av=\lambda v\): sólo algunas direcciones, o ninguna real si hay giro & una rotación de
\(\R^3\) tiene exactamente una: el eje & dirección que cae sobre sí misma, sólo estirada o
invertida\\
\addlinespace
Derivada & \(f'(x)\): un número, la pendiente & gradiente: una forma lineal con dos
números; jacobiano \(2\times2\) si la función es vectorial & gradiente con \(n\) números;
jacobiano \(m\times n\) & la mejor aproximación lineal: recta tangente, plano tangente\\
\addlinespace
Recta & \(a\,x=b\): un punto & forma lineal \(n\cdot x=c\): una recta; al variar \(c\),
rectas paralelas & \(n\cdot x=c\): un plano; en \(\R^n\), hiperplano & el conjunto de
nivel de una regla graduada (sección~\ref{sec:formas-lineales})\\
\bottomrule
\end{tabularx}
}

\begin{idea}[Cómo usar la tabla]
Ante una fórmula nueva con matrices, escribirla primero con un solo número. Si en \(\R\)
la fórmula es obvia, la versión matricial dice lo mismo dirección por dirección; lo único
que agregan las dimensiones es que ahora las direcciones pueden mezclarse, y por eso
aparecen el orden de composición, los autovectores y el determinante.
\end{idea}

\section{Del álgebra a la geometría: qué pregunta cada operación}\label{sec:algebra-geometria}

\textbf{Desarrollo.} El hilo de la materia es siempre el mismo: una operación algebraica,
en apariencia abstracta, resulta ser exactamente la pregunta geométrica que necesitábamos.
Conviene tener el mapa antes de las cuentas.

\begin{center}
\begin{tabularx}{\textwidth}{>{\bfseries}p{34mm} X X}
\toprule
Operación & Qué calcula & Qué pregunta geométrica responde\\
\midrule
Producto interno \(u\cdot v\) & un número & ¿son ortogonales? ¿cuánto se proyecta uno
sobre otro? En homogéneas: ¿el punto está en la recta?\\
Producto vectorial \(u\times v\) & un vector normal a ambos & ¿cuál es la recta que une dos
puntos? ¿cuál es el punto donde se cortan dos rectas? ¿qué área orientada forman?\\
Determinante \(\det[u\ v\ w]\) & un número & ¿están alineados tres puntos? ¿son
linealmente dependientes tres vectores? ¿qué volumen orientado forman?\\
Producto exterior \(u\wedge v\) & un bivector (área orientada) & lo mismo que el
vectorial, pero en cualquier dimensión\\
Núcleo de \(A\) & un subespacio & ¿qué vectores aplasta \(A\) a cero? Si hay alguno,
\(\det A=0\)\\
Autovectores de \(A\) & vectores con \(Av=\lambda v\) & ¿qué direcciones deja quietas
\(A\)? En homogéneas: ¿qué puntos deja fijos la transformación?\\
\bottomrule
\end{tabularx}
\end{center}

\begin{idea}[La cadena]
incidencia \(\to\) unión de puntos \(\to\) alineación \(\to\) determinante \(\to\)
núcleo \(\to\) autovectores \(\to\) puntos fijos. Cada eslabón se deduce del anterior;
es lo que desarrolla la segunda lectura.
\end{idea}

\begin{idea}[En palabras del profesor]
``Traemos a la mesa a nuestro `detector de colinealidad': el producto cuña, y lo usamos
primero para construir un polinomio cuyas raíces son los autovalores, y luego para hallar
los autovectores como el subespacio nulo de unas matrices que armamos'' (resumen de clase 7, cursada 2025).
\end{idea}

\section{Transformaciones proyectivas: conservar rectas, nada más}\label{sec:proyectivas-primera}

\textbf{Dictado}: \emph{``Transformaciones proyectivas: mantienen rectas las rectas
\(\Rightarrow\) quedan definidas por un cuadrilátero''}.

\textbf{Desarrollo.} La imagen mental es una fotografía. Una cámara inclinada sobre un
terreno plano convierte un lote rectangular en un trapecio: los alambrados paralelos
convergen hacia un punto de fuga, las distancias no se conservan en ninguna proporción fija,
y sin embargo toda recta del terreno sigue siendo una recta en la foto. Esa es la única
propiedad que sobrevive, y es la que define a la familia.

\begin{idea}[Qué agrega la perspectiva]
Las afinidades conservan rectas \emph{y paralelismo}. Las transformaciones proyectivas
conservan rectas \emph{solamente}: dos paralelas pueden dejar de serlo y cortarse en un punto
de fuga; un punto ordinario puede irse al infinito y un punto ideal puede volverse ordinario.
Las afinidades son un caso particular de las proyectivas.
\end{idea}

En coordenadas homogéneas la diferencia es mínima y muy visible. Una afinidad tiene matriz
con última fila \((0,0,1)\), de modo que \(W\) siempre vale \(1\). Una transformación
proyectiva tiene una matriz \(3\times3\) invertible \emph{cualquiera}:
\[
\begin{pmatrix}X\\Y\\W\end{pmatrix}=H\begin{pmatrix}x\\y\\1\end{pmatrix},
\qquad
(x',y')=\left(\frac{X}{W},\frac{Y}{W}\right).
\]
Ahora \(W=h_{31}x+h_{32}y+h_{33}\) depende del punto, y la división por \(W\) es
exactamente la perspectiva: lo que ``está lejos'' se achica. Los puntos con \(W=0\) forman
una recta que se va al infinito, el horizonte de la foto.

\paragraph{Por qué un cuadrilátero.}
La matriz tiene nueve entradas, pero \(H\) y \(\lambda H\) representan la misma
transformación (el factor se cancela al dividir por \(W\)). Quedan ocho números
independientes. Cada correspondencia ``el punto \(p\) va al punto \(p'\)'' aporta dos
ecuaciones, una por coordenada. Cuatro puntos, sin tres alineados, dan ocho ecuaciones: la
transformación queda determinada. Con tres puntos falta decidir ``cuánta perspectiva hay'';
con cinco puntos, en general, no hay solución, y el quinto va a donde la transformación lo
mande.

\paragraph{Puntos fijos.}
Un punto es fijo cuando \(Hp=\lambda p\): un autovector de \(H\). El polinomio característico
tiene grado \(3\), y un polinomio real de grado impar siempre tiene una raíz real. Por lo tanto
\textbf{toda transformación proyectiva del plano tiene al menos un punto fijo}, que puede ser
un punto ordinario o un punto ideal (una dirección fija, como en una traslación).

\begin{idea}[En palabras del profesor]
``Estos espacios invariantes ya habían aparecido cuando usábamos las coordenadas
homogéneas para analizar transformaciones afines o proyectivas. En ese caso los subespacios
invariantes de la transformación lineal de \(\R^3\) correspondían a los puntos fijos en
\(\mathbb{P}\R^2\)'' (resumen de clase 6, cursada 2025).
\end{idea}

